\documentclass[11pt, a4paper]{article}
\usepackage[utf8]{inputenc}
\usepackage[T1]{fontenc}
\usepackage[margin=1in]{geometry}
\usepackage[dvipsnames]{xcolor}
\usepackage{titlesec}
\usepackage{parskip}
\usepackage{enumitem}

% --- Custom Styling ---
\titleformat{\section}
{\normalfont\Large\bfseries\color{NavyBlue}}{\thesection}{1em}{}

\titleformat{\subsection}
{\normalfont\large\bfseries\color{Black}}{\thesubsection}{1em}{}

% Define color commands for consistent styling
\newcommand{\bad}[1]{\textbf{\textcolor{BrickRed}{BAD: #1}}}
\newcommand{\good}[1]{\textbf{\textcolor{ForestGreen}{GOOD: #1}}}

% --- Title Information ---
\title{\textbf{The Daily Gamecock \textit{Opinion} \\ Headline Style Guide}}
\author{J.C. Vaught, Opinion Editor. Fall 2025}
\date{}

\begin{document}

\maketitle

\section{The Core Philosophy}

\textbf{We write arguments, not breaking news.}

A news headline tells the reader \textit{what} happened. An opinion headline tells the reader \textit{why it matters} and \textit{who is responsible}, and often, \textit{what should be done about it}. We do not \textit{only} report on the existence of a problem; we pass judgment on it.

\subsection{Principle 1: The ``No Colons'' Rule}
Colons are for lab reports, textbooks, and news features. They create a ``Label: Explanation'' dynamic that is boring and passive. Opinion headlines ought to be fluid, declarative sentences that can stand on their own. It also looks weird to have 2 colons in a title; especially since all of our columns need to start with [Column: title] or [Guest Column: title].

\textbf{EXAMPLE 1 -- Dining:} A writer argues that the new meal plan is too expensive and food quality has dropped. 
   
    \bad{Column: Dining Services: The new meal plan is a failure} 
    
    \good{Column: The new dining contract exploits a captive freshman class} 
    
    The bad column was reasonably effective, but it was missing some key elements -- namely, WHY the plan was a failure. Additionally, it is quite clear to any competent reader that the meal plan is being run by Dining Services; repeating that information is a waste of time for both the writer and the reader.
   
    The revised column title also answers WHO was affected -- in this case, the freshman class.
   
    The new title also states WHAT happened. While not required, this is helpful for readers so they do not waste time clicking on an article that is useless or uninteresting to them.
    
    The original piece just tells me that the author has an opinion about the new meal plan. The new piece tells me that the author thinks the new contract (the what) is hurting the freshman class (the who).

\newpage
EXAMPLE 2: \textbf{Context (Sports).} The football team lost largely because the defense could not stop the run.

\bad{Column: Gamecock Football: Defense was the problem on Saturday}

\good{Column: A porous defense sabotages our offensive efforts}

The bad column tells the reader which side of the ball struggled, but it does not say how or why that struggle mattered. It simply labels ``defense'' as the problem without making an argument about consequences.

The revised headline turns that label into a claim: the defense did not just play poorly; it \textit{sabotaged} the offense. That verb signals blame and implies a causal chain from missed tackles and blown assignments to stalled drives and a loss on the scoreboard.

The bad headline also wastes space on information that is either redundant or unhelpfully vague. Because we operate on a multi-week writing and publishing schedule, saying that something happened ``on Saturday'' is imprecise; readers may not know which Saturday is meant, and the phrase adds no argumentative value. Likewise, including ``Gamecock Football'' is unnecessary context at a South Carolina–focused publication, where readers can safely assume that any reference to ``football'' is about Gamecock football. Those words consume headline real estate without sharpening the claim.

Even if the writer wanted to stick with a straightforward descriptive thesis, a line like ``Gamecock defense was the problem'' would communicate the same idea without the extra colon and label–explanation structure. The stronger revised headline goes further by hinting at the stakes. By pairing ``porous defense'' with ``our offensive efforts,'' it frames the story as one of wasted opportunities, not just a generic bad performance, and it gives the writer a clear argumentative lane to show how specific defensive failures undercut what the offense did well.


EXAMPLE 3:  \textbf{Context (Academics):} The library is cutting its 24-hour schedule during finals week. 

    \bad{Column: Library Hours: Why they need to change back} 
    
    \good{Column: Reduced library study hours has sabotaged our academic success}
    
    The bad headline promises that the writer will explain ``why they need to change back,'' but it does not tell the reader who is harmed or what exactly is at stake. It reads more like a teaser for a list of complaints than a thesis about a specific failure.
    
The revised headline directly names both the policy and its impact. ``Reduced library study hours'' identifies the concrete change, and ``have sabotaged our academic success'' tells the reader that this is not a minor inconvenience but a serious barrier to doing well in classes.

By explicitly linking administrative decisions about scheduling to outcomes like grades, stress, and equity, the new headline sets up an argument rather than a mere preference. It invites the writer to prove that this policy makes it harder for students -- especially those with jobs or caregiving responsibilities -- to succeed academically, instead of just insisting that longer hours would be more convenient.


\subsection{Principle 2: The ``No False Urgency'' Rule}
Breaking news relies on words like ``Just,'' ``Now,'' ``Recently,'' or ``Yesterday.'' Opinion writing is timeless. If your argument is strong, it does not need a timestamp to make it relevant. Focus on the \textit{implication}, not the \textit{timing}.

\begin{itemize}
    \item \textbf{Context (Student Gov):} Student Government passed a bill yesterday to fund a new aesthetic garden instead of student orgs. \\
    \bad{Student Senate just passed a bill to build a garden instead of academics} \\
    \good{Funding aesthetic gardens over student scholarships is a gross misuse of fees}

    \item \textbf{Context (Tuition):} The Board of Trustees announced a tuition hike this morning \\
    \bad{The Board just voted to raise tuition again --- third year in a row} \\
    \good{Tuition hikes demand an immediate freeze on administrative salaries}

    \item \textbf{Context (Safety):} A robbery occurred last night at an off-campus apartment complex. \\
    \bad{A robbery just happened at The Hub last night} \\
    \good{Recurring crime in off-campus housing exposes a failure in city zoning oversight}
\end{itemize}

\subsection{Principle 3: The ``No Explainers'' Rule}
Never promise to explain something to the reader (``Why/How X happened''). That is the job of a reporter. As a columnist, your job is to deliver the verdict immediately. Turn the ``Why'' into a ``Result.''

\begin{itemize}
    \item \textbf{Context (Athletics):} The Head Coach fired the Offensive Coordinator after a losing streak. \\
    \bad{Why the Head Coach decided to fire the Offensive Coordinator} \\
    \good{Firing the coordinator is a desperate attempt to save a failing season}

    \item \textbf{Context (Greek Life):} Fewer students are rushing sororities this year compared to last year. \\
    \bad{Why sorority recruitment numbers are dropping} \\
    \good{Declining recruitment numbers signal a cultural shift away from Greek Life}

    \item \textbf{Context (Transportation):} The university changed the bus routes, and now they do not go to the engineering building. \\
    \bad{How the new bus routes affect engineering students} \\
    \good{The new bus schedule effectively isolates engineering majors from main campus}
\end{itemize}

\subsection{Principle 4: The ``Opinion as Fact'' Rule (No Generalizations)}
New writers often try to hide their opinion by saying ``Students feel...'' or ``Many think...'' This is weak. Do not report on the students' feelings; report on the \textit{validity} of those feelings. State the opinion as an objective truth.

\begin{itemize}
    \item \textbf{Context (Parking):} Commuters are angry because the university oversold parking permits. \\
    \bad{Students are frustrated and angry about the lack of parking spots} \\
    \good{Overselling permits is a predatory practice that steals from commuters}

    \item \textbf{Context (Housing):} Students are scared because they found mold in the Capstone dorms. \\
    \bad{Residents are worried about mold in their rooms} \\
    \good{Unchecked mold growth in residence halls is public health negligence}

    \item \textbf{Context (Mental Health):} Students feel that the wait times for counseling services are too long. \\
    \bad{Students think counseling services are too slow} \\
    \good{Month-long wait times for counseling render the service useless during a crisis}
\end{itemize}

\subsection{Principle 5: The ``No Names'' Rule (Focus on Roles)}
Unless the person is iconic (The President, The Head Coach), using a specific name in a headline often alienates readers who do not know them. Focus on the \textit{role}, the \textit{action}, or the \textit{impact}. The person is temporary; the structural issue is permanent.

\begin{itemize}
    \item \textbf{Context (Student Gov):} Senator John Doe voted against a bill providing free scantrons. \\
    \bad{John Doe voting no on scantrons was a mistake} \\
    \good{Denying academic supplies during finals week signals a disconnect from student needs}

    \item \textbf{Context (Sports):} The starting quarterback keeps throwing interceptions in the fourth quarter. \\
    \bad{Smith needs to stop throwing interceptions} \\
    \good{Reckless quarterback play is the single greatest liability to this season}

    \item \textbf{Context (Admin):} Provost Smith decided to change the grading scale to remove minuses. \\
    \bad{Provost Smith made a bad call on the grading scale} \\
    \good{Arbitrary changes to the grading scale undermine academic integrity}
\end{itemize}

\subsection{Principle 6: The ``Active Blame'' Rule}
Passive voice protects the guilty. Never say ``Mistakes were made'' or ``Problems exist.'' Identify exactly \textit{who} is doing the action. If a system is broken, name the system.

\begin{itemize}
    \item \textbf{Context (Facilities):} The elevators in a dorm have been broken for three weeks. \\
    \bad{Broken elevators are a major inconvenience} \\
    \good{Housing neglect forces residents to live in inaccessible conditions}

    \item \textbf{Context (Technology):} The Wi-Fi in the library crashes every day at noon. \\
    \bad{The Wi-Fi in the library is really bad} \\
    \good{Infrastructure failures in the library prevent students from completing basic coursework}

    \item \textbf{Context (Events):} The ticket lottery system crashed, and nobody got tickets. \\
    \bad{Mistakes were made with the student ticket lottery} \\
    \good{Athletics failed to prepare for predicted student demand}
\end{itemize}

\subsection{Principle 7: The ``So What?'' Rule (Future Consequence)}
A complaint looks backward (``This is bad''). An argument looks forward (``This will cause X''). The best headlines predict the consequence of the status quo.

\begin{itemize}
    \item \textbf{Context (Academics):} The university is doubling class sizes for English 101 to save money. \\
    \bad{English classes are getting too big for students to learn} \\
    \good{Doubling class sizes ensures the erosion of the liberal arts education}

    \item \textbf{Context (Diversity):} The state legislature is cutting the budget for the diversity office. \\
    \bad{Cutting the diversity budget is bad for the school} \\
    \good{Defunding diversity programs ensures a homogenized and unprepared student body}

    \item \textbf{Context (City Life):} A beloved local coffee shop is being replaced by a Starbucks. \\
    \bad{It is sad that Cool Beans is closing down} \\
    \good{Allowing corporate chains to dominate downtown erases the city's cultural identity}
\end{itemize}

\newpage

\section{Approved 'analysis'-style Headline Structures}

These are the ``plug-and-play'' formulas for writing strong headlines. When you are stuck, see which of these four templates fits your argument best.

Note that if you are writing a `call-to-action' piece or a `activism'-esque piece, then you should skip to section 3 where we go over this in more detail. This section is more for analysis-style or reflection columns and opinions

\subsection{Structure 1: The ``Dead/Over'' Declaration}
This structure declares the end of an era, a tradition, or a strategy. It is definitive and bold. It does not ask \textit{if} something is over; it asserts that it \textit{is}.

\begin{itemize}
    \item \textbf{Context (Tradition):} A columnist argues that mandatory attendance policies are outdated in the age of digital learning. \\
    \bad{Mandatory Attendance: Why professors should stop using it} \\
    \good{The era of mandatory attendance is dead and that is progress}

    \item \textbf{Context (Athletics):} The Athletic Department keeps trying to sell season tickets by highlighting championships from the 1980s, but students don't care. \\
    \bad{Athletics needs to stop focusing on the past to sell tickets} \\
    \good{Relying on alumni nostalgia to sell student tickets is a dead strategy}

    \item \textbf{Context (Campus Culture):} Students are moving off-campus earlier and earlier, killing the traditional ``dorm life'' culture. \\
    \bad{Dorm culture isn't what it used to be} \\
    \good{The concept of a unified `campus life' is dead to a commuter-heavy class}
\end{itemize}

\subsection{Structure 2: The Comparative Value (``X matters more than Y'')}
This structure forces a choice between two values. It acknowledges that both might be good, but argues that one is \textit{essential}. It creates conflict and stakes.

\begin{itemize}
    \item \textbf{Context (Student Fees):} Student Government wants to renovate the student union lobby (aesthetic) instead of lowering the price of meal plans (affordability). \\
    \bad{We need lower food prices more than a new lobby} \\
    \good{Affordable dining matters more than aesthetic renovations to meal plan holders}

    \item \textbf{Context (Elections):} A candidate with a perfect resume lost to a candidate who actually talked to students. \\
    \bad{Students prefer candidates who are nice over candidates with resumes} \\
    \good{Authentic connection matters more than a resume to this student body}

    \item \textbf{Context (Infrastructure):} The university is prioritizing building more parking garages over making crosswalks safer for pedestrians. \\
    \bad{Safety is more important than parking spots} \\
    \good{Pedestrian safety matters more than convenient parking spots}
\end{itemize}

\subsection{Structure 3: The Action/Consequence (``A acts as/leads to B'')}
This is the most ``scientific'' of the opinion headlines. It posits that \textbf{Action A} inevitably causes \textbf{Result B}. Use strong, active verbs like \textit{ensures, guarantees, signals,} or \textit{acts as}.

\begin{itemize}
    \item \textbf{Context (Housing):} The city refuses to put a cap on rent prices for student apartments, so poor students are dropping out. \\
    \bad{High rent prices are causing students to leave school} \\
    \good{Ignoring the housing crisis acts as a silent eviction for low-income students}

    \item \textbf{Context (Tuition):} The university raised tuition but the buildings are still falling apart. \\
    \bad{Tuition is going up but facilities aren't getting better} \\
    \good{Raising tuition without improving facilities is a breach of trust}

    \item \textbf{Context (Gentrification):} The university is tearing down cheap local bars to build high-end luxury condos. \\
    \bad{Building luxury condos is hurting the local community} \\
    \good{Prioritizing luxury apartments ensures the gentrification of the downtown district}
\end{itemize}

\subsection{Structure 4: The Inverse/Reframe}
This structure takes a common narrative (what ``everyone'' says) and flips it on its head. It shifts the blame from the victim to the system.

\begin{itemize}
    \item \textbf{Context (Scooters):} Everyone hates electric scooters on sidewalks, but the columnist argues they only ride there because there are no bike lanes. \\
    \bad{Scooters aren't the problem, bike lanes are} \\
    \good{The demonization of electric scooters is a lazy alternative to fixing bike lanes}

    \item \textbf{Context (Attendance):} Professors complain that students aren't coming to class, but the columnist argues the classes are boring and useless. \\
    \bad{It's not students' fault they don't go to class} \\
    \good{Blaming students for low attendance ignores the failure of the lecture format}

    \item \textbf{Context (Activism):} The university deleted the designated ``Free Speech Zone,'' claiming it was to ``protect students,'' but it actually silences them. \\
    \bad{Removing the Free Speech Zone is bad for protests} \\
    \good{The removal of the free speech zone is the first step in sanitizing campus culture}
\end{itemize}

\newpage

\section{Activism \& Demands}

\textbf{We do not beg; we demand.}

Weak headlines ask for change (``The university should consider X''). Strong headlines mandate change by highlighting the consequences of inaction. Avoid preachy language like ``Students need to...'' or ``You should...''

\subsection{Structure 1: The ``It Is Time'' (The Ultimatum)}
This structure is the strongest way to call for a total overhaul. It implies that patience has run out. Use it when a system is fundamentally broken and cannot be fixed with small tweaks.

\begin{itemize}
    \item \textbf{Context (Policy):} The current attendance policy penalizes students for being sick, even with a doctor's note. \\
    \bad{The attendance policy needs to be more lenient} \\
    \good{It is time to dismantle the outdated mandatory attendance policy}

    \item \textbf{Context (Business Ties):} The current dining vendor has failed health inspections three times in a row. \\
    \bad{We should probably find a new dining company} \\
    \good{It is time to sever ties with the current dining vendor contract}

    \item \textbf{Context (Bureaucracy):} Student Activity Fees are being spent on parties for the executive officers, and no one knows where the rest of the money is. \\
    \bad{Student Gov needs to show us where the money is going} \\
    \good{It is time to audit the mishandled allocation of student activity fees}
\end{itemize}

\subsection{Structure 2: The Mandate (``Must,'' not ``Should'')}
``Should'' is a suggestion; ``Must'' or ``Need'' is a requirement. When targeting a specific powerful entity (The Board, City Council, The President), use ``Must'' to sound authoritative. Connect the action to a specific result.

\begin{itemize}
    \item \textbf{Context (Tuition):} Tuition keeps going up, and students are dropping out because they can't afford it. \\
    \bad{The Board of Trustees should stop raising tuition} \\
    \good{The Board of Trustees must freeze tuition hikes to protect retention rates}

    \item \textbf{Context (City Planning):} Developers are building luxury student housing that no actual student can afford. \\
    \bad{The city should make developers build cheaper apartments} \\
    \good{City Council must force developers to include affordable housing caps}

    \item \textbf{Context (Housing Monopoly):} One company owns all the off-campus apartments and is price-gouging students. \\
    \bad{The administration needs to say something about the apartment prices} \\
    \good{The Administration must address the off-campus housing monopoly}
\end{itemize}

\subsection{Structure 3: The ``A, Not B'' (The Responsibility Shift)}
This structure is excellent for reframing a debate. Often, institutions blame students for problems (safety, trash, noise). This structure pushes the blame back up the ladder to the people with power.

\begin{itemize}
    \item \textbf{Context (Safety):} The university tells students to ``walk in groups'' to stay safe, rather than hiring more security or fixing the streetlights. \\
    \bad{Students shouldn't be the only ones worrying about safety} \\
    \good{The Administration, not the student body, must shoulder the cost of security}

    \item \textbf{Context (Labor):} The university relies on underpaid graduate students to teach classes while tenured professors do research. \\
    \bad{Grad students need more help from professors} \\
    \good{Tenured faculty, not graduate assistants, must lead the charge for fair wages}

    \item \textbf{Context (Infrastructure):} Downtown roads are crumbling because of the heavy construction equipment used to build student condos. \\
    \bad{We shouldn't have to pay for the roads the developers broke} \\
    \good{Developers, not residents, need to pay for the downtown infrastructure strain}
\end{itemize}

\end{document}